\chapter{PROCESO ITERATIVO DE DISEÑO}
\label{anx:iteraciones}

Este anexo registra los cambios técnicos y paramétricos entre las revisiones del diseño, por subsistema, y respalda la Tabla~\ref{tab:iteraciones-resumen}; la Tabla~\ref{tab:anx-iteraciones} reúne el registro. Una revisión solo sustituye a la anterior en el aspecto que se comprobó; cada resultado se asigna a la geometría y al caso que lo produjeron. Las fechas corresponden a 2026.

\begingroup
\small\singlespacing
\setlength{\tabcolsep}{4pt}
\begin{longtable}{@{}>{\raggedright\arraybackslash}p{2.6cm}>{\raggedright\arraybackslash}p{\dimexpr(\textwidth-2.6cm-2.4cm-8\tabcolsep)*55/100\relax}>{\raggedright\arraybackslash}p{\dimexpr(\textwidth-2.6cm-2.4cm-8\tabcolsep)*45/100\relax}>{\raggedright\arraybackslash}p{2.4cm}@{}}
\caption{Registro de iteraciones por subsistema}\label{tab:anx-iteraciones}\\
\toprule
\textbf{Revisión} & \textbf{Cambio} & \textbf{Motivo} & \textbf{Estado}\\
\midrule\endfirsthead
\caption[]{Registro de iteraciones por subsistema (continuación)}\\
\toprule
\textbf{Revisión} & \textbf{Cambio} & \textbf{Motivo} & \textbf{Estado}\\
\midrule\endhead
\midrule\multicolumn{4}{r}{\textit{Continúa en la página siguiente.}}\\\endfoot
\bottomrule\endlastfoot
\multicolumn{4}{@{}l}{\textit{Estructura y mesa}}\\
Disposición inicial (julio) & Mesa abierta conceptual & Arranque del proyecto & Histórica\\
Modelos de cabina (agosto) & Mesa circular con aros, nervios y brazos de rueda en voladizo & Evolución de la arquitectura de la cabina & Superada\\
Modelos de línea (28--30/08) & Bastidores de tres, cuatro y cinco apoyos sectoriales con el mismo perfil & Comparar topologías para el análisis estructural & Superada\\
Análisis estructural (31/08--04/09) & Bastidor de tres apoyos en W6$\times$9: D/C de 0,28 & Cribado estructural & Ejecutado\\
Revisiones mecánicas 1--3 (05--06/09) & Cabina y puerta, rieles A45, recorrido de 3320 mm, masa de 3500 kg, ruedas SPKVS 200K; horquilla de 68 a 80 mm; reacciones en los apoyos por equilibrio & Desarrollo mecánico y revisión independiente & Superadas\\
Revisión 4 (07--08/09) & Bloqueo unido al bastidor; centrador retraído; eje de rueda de 35 mm & Encargo de revisión mecánica & Superada\\
Revisión 5 (08/09) & Colocación $\leq5$ mm; topes y manillar conectados; envolvente del riel & Cierre de los hallazgos de la revisión 4 & Vigente (mesa)\\
Integración (24/09 en adelante) & Soportes de riel reubicados (12 por riel); piso de drenaje fijo con pendiente, sumidero y cesta; estados de carga y lavado & Choques del centrador con la jamba, de los soportes con el muro y del drenaje con los rieles & Candidato\\
Travesaño adicional (29/09) & Propuesto y retirado & Mejoraba la deformación, pero quedaba bajo la lanza y dependía de uniones a momento & Retirado\\
Rueda y riel (27--29/10) & Fuerza lateral de guiado de 6,51 kN; retención axial de las ruedas; se compararon las ruedas de 200 y 250 mm & La versión anterior comparaba la carga última con la capacidad de rodadura & Se mantiene la de 200 mm\\
Traslado (29/10) & Doble cadena bajo la mesa con volante en lugar del empuje manual & Empuje de 325 a 837 N y frenado de 565 N & Vigente\\
Uniones del bastidor (02/11) & Alas a tope con penetración completa, alma con filete de 3 mm, rigidizador de 6 mm y tapas en las esquinas & Las uniones estaban especificadas pero no diseñadas & Vigente\\
Pedestales del riel (02--03/11) & Pedestal de inoxidable 304 con 4 anclajes adhesivos M16 & Levantamiento de 3,78 kN, guiado y tensor de las cadenas & Vigente\\
Topes de retención (03/11) & Cartelas a 60$^\circ$ y 300$^\circ$ y chapa de asiento a 190$^\circ$ & La base supuesta del tope no existía en el CAD: cordón con D/C 3,97 & Vigente\\
\multicolumn{4}{@{}l}{\textit{Centradores}}\\
Módulos de carro radial (revisión 5) & Centradores con carro y guía & Centrado de la revisión 5 & Superados: atravesaban perfiles e invadían el rodamiento\\
Selección de conceptos & Carro radial, brazo pivotante y leva excéntrica; puntajes de 2,675, 3,250 y 3,075 & Diez condiciones excluyentes y matriz ponderada & Brazo pivotante seleccionado\\
Versiones 5 a 5.7.1 & Desarrollo del brazo pivotante, con dos referencias y un empujador & Modelado, barridos y elementos finitos & Superadas\\
Versión 5.7.2b & Amarre transversal de las referencias con dos pernos M12 & Flexibilidad de la mordaza & Superada\\
Versión 5.7.3 & Brazo del empujador de 67,15 mm (+2 mm de recorrido) & Captura de 4,69 mm, menor que 5 mm, en el peor caso & Superada\\
Versión 5.7.4 (20/10) & Placa del soporte del empujador de 20 a 28 mm, con soldaduras a tope de penetración completa & Tensiones de 190--246 MPa en la placa de 20 mm & Vigente\\
Amarre del empujador (27/10) & Dos pernos pasantes M12 por el alma de la viga longitudinal & El apriete anterior no tenía un camino de reacción cerrado & Vigente\\
Centrado por fuerza (02/11) & Ajuste H7/g6 y límite de 1,95 kN a la fuerza final; sin cono de indexado & Error de 6,10 mm en la envolvente; el 70\,\% era deformación elástica & Vigente\\
\multicolumn{4}{@{}l}{\textit{Aspersión}}\\
Modelos de cabina (agosto) & Distribuidor en U, brazos dobles, brazo en L y cabezal sobre monorriel & Cobertura geométrica por trazado de rayos & Superados\\
Auditoría hidráulica (01/09) & Arquitecturas del cabezal sin seleccionar & Faltaban datos de proceso & Superada\\
U giratoria (12/09) & La U superior giratoria sustituye al cabezal sobre monorriel & Nuevo planteamiento del aplicador & Vigente\\
Caso base (16/09) & 19 boquillas, 1,7 rpm, agua a 62~$^\circ$C, 174 L/min & Memoria hidráulica del caso base & Superado por la red de 22\\
Lanza corta (18/09) & Lanza a 25$^\circ$ en lugar de 28$^\circ$; dos salidas inferiores exteriores; 22 boquillas; traslado solo con el aplicador a 0$^\circ$ & Con 28$^\circ$ quedaban 90 de 18\,720 puntos sin acceso & Vigente, salvo las salidas inferiores\\
Cierre preliminar (18/09) & Manguera en lazo, toma fija reubicada, válvulas, purgas y seguros; boquilla comercial candidata & Cierre de ingeniería preliminar del aplicador & Vigente\\
Red de 22 boquillas & 202,59 L/min y 52,56 m; toma crítica en la lanza & Geometría de la lanza corta & Superada por la de 24\\
Par inferior de abanico estrecho (08/10) & Boquilla inferior exterior a 15$^\circ$ y una adicional a 10,7$^\circ$ en cada brazo: 24 boquillas, 221,07 L/min y 53,39 m & La cara inferior recibía cuatro veces menos dosis; la boquilla anterior mojaba la base del diámetro exterior & Vigente\\
U desmontable (02/11) & Dos juntas bridadas DN50 a 300 mm del eje, sin válvulas de ramal; brazos iguales & Las válvulas de ramal no tenían función documentada & Vigente\\
Rotor, eje y cubo (02/11) & Masa medida en el CAD: 231 kg; cubo soldado al eje y mecanizado con él & La envolvente de 180 kg no cubría el rotor real & Vigente\\
Tolerancias de las coplas (03/11) & $\pm2$ mm y $\pm1^\circ$ en las coplas; $\pm1$ mm y $\pm0{,}5^\circ$ en el bloque inferior & Trazado de rayos con 240 perturbaciones & Vigente\\
Soporte superior (27/10 y 03/11) & Marco propio dentro de los muros; postes SHS80$\times$4 y vigas RHS150$\times$100$\times$3 & La cabina no basta como estructura; el criterio anterior no era funcional & Vigente: de 360 a 247 kg\\
\multicolumn{4}{@{}l}{\textit{Simulación del mojado}}\\
Planificación (julio a 18/09) & De una boquilla aislada a un modelo de gotas de las 6 boquillas del agujero & Definir el fenómeno relevante & Superada\\
Agujero en OpenFOAM (19--21/09) & Dos ramas de cálculo; la oficial, con acoplamiento en dos vías y 430\,294 celdas & Primer modelo del mojado & Antecedente\\
Caso puente en Fluent (26--30/09) & Agujero repetido en Fluent: $-0{,}27$\,\% en dosis media & Ampliar el dominio a 22 boquillas dentro del límite de la licencia & Ejecutado\\
Caso Base de 22 boquillas (30/09--02/10) & Cuatro superficies con los obstáculos del soporte; verificación de malla, paso de tiempo, dominio y parcelas & Configuración vigente en esa fecha & Ejecutado\\
Variantes del diámetro exterior & Lateral 1 inclinada $-4^\circ$; segunda boquilla superior 15 mm hacia fuera & Franja seca de 2 mm en el diámetro exterior & Lateral inclinada incorporada (26/10); superior desplazada no incorporada\\
Configuración propuesta (25--26/10) & 24 boquillas con los dos brazos explícitos; tres mallas con Courant de gota constante y dos números de parcelas & Comprobar el par inferior con el CFD & Ejecutado\\
\end{longtable}
\endgroup

La Tabla~\ref{tab:comparacion-apoyos} compara los bastidores de tres, cuatro y cinco apoyos con la misma carga externa total, repartida por igual entre sus apoyos. Se mantuvo el de tres apoyos por masa y simplicidad, a cambio de una flecha 63\,\% mayor que la del de cinco, todavía inferior a los límites en evaluación.

\begin{table}[htbp]
\centering\small\singlespacing
\caption{Resultados de las alternativas de bastidor}
\label{tab:comparacion-apoyos}
\begin{tabular}{lrrr}
\toprule
\textbf{Magnitud} & \textbf{Tres apoyos} & \textbf{Cuatro apoyos} & \textbf{Cinco apoyos}\\
\midrule
Carga externa por apoyo (N) & 9948 & 7461 & 5968,8\\
Elementos de barra & 20 & 24 & 28\\
Masa de perfiles (kg) & 168,1 & 188,5 & 180,8\\
Máximo $|U_3|$ nominal (mm) & 1,353 & 1,441 & 0,953\\
Flecha relativa nodal (mm) & 1,236 & 1,415 & 0,757\\
Desnivel entre apoyos (mm) & 0,294 & 0,410 & 0,220\\
D/C publicado & 0,277 & 0,278* & 0,175*\\
\bottomrule
\end{tabular}
\par\smallskip\footnotesize
*Combinación resistente cuya equivalencia con la del bastidor de tres apoyos no está demostrada; no se usa para ordenar las alternativas.
\end{table}
